\section{[INTERNAL] Entity Identity Across the Corpus}
\label{sec:appendix-internal-entity-iris}
%% ====================================================================
%% INTERNAL appendix --- visible only when compiling draft.tex.
%% Captures the design discussion (Daniel + reviewer agent, 2026-04-27)
%% on how to keep entity IRIs consistent as the corpus grows.
%% Do NOT include in submission.
%% ====================================================================

\begin{center}
\fbox{\parbox{0.92\textwidth}{\centering\small
\textbf{INTERNAL --- not part of the submission.}\\
This appendix records design decisions on entity-IRI hygiene
in the per-paper KG corpus and is included in the
\texttt{draft.tex} build only.}}
\end{center}

\bigskip

The corpus is encoded one paper at a time, by an agent that does
not see other papers' encodings during its session. This pattern
naturally exposes two failure modes: \emph{splits} (the same
real-world entity ending up with two distinct IRIs across two
papers) and \emph{collisions} (the same IRI denoting two distinct
entities in two papers). This appendix documents the conventions
we adopt to keep both rare.

\subsection{The two failure modes}

\paragraph{Split: same entity, two IRIs.}
Two papers both use the PBE exchange--correlation functional, but
one encoding mints \texttt{ex:xc-pbe-foo2024} and another
\texttt{ex:pbe-functional-bar2025}. SPARQL exact-match queries see
two distinct entities; OWL reasoning has no basis to merge them;
aggregate counts (``how many papers used PBE?'') fragment.

\paragraph{Collision: same IRI, different entities.}
Two papers both use the local IRI \texttt{ex:dft-settings} but for
different DFT settings (one means VASP+PBE+400~eV, the other means
GPAW+SCAN+500~eV). Cross-corpus merges then produce contradictory
triples on the same subject.

\subsection{Mitigations}

\paragraph{Collision is solved by namespace discipline.}
Every paper-local IRI carries the paper-id as a suffix:
\texttt{ex:dft-settings-kumar2025-c15} rather than
\texttt{ex:dft-settings}. The executor agent already follows this
convention; we promote it to a hard rule. When the corpus is
loaded into Concon as one named graph per paper, even an accidental
collision is contained to a single graph.

\paragraph{Split is solved by a controlled vocabulary plus external
identifiers.}
Concepts that recur across papers (XC functionals, pseudopotential
types, wave-function methods, software packages, MLIP method
families, common dispersion corrections) are pre-minted as named
individuals in a shared vocabulary file
\texttt{artifacts/kg/mlips-vocab.ttl}. The encoding instruction
becomes ``if PBE is mentioned, use \texttt{mlips:PBE}''---the agent
selects from a list rather than minting locally. This is the
pattern already used for $\DatasetProvenanceC$
(\texttt{mlips:Published}, \texttt{mlips:InHouse}, \dots),
$\MetricTypeC$ (\texttt{mlips:RMSE}, \texttt{mlips:MAE}), and
$\SamplingStrategyC$.

External identifiers act as a second layer of alignment:
$\dlRole{sameAsWikidata}$ for materials, software, and persons;
DOIs for papers; CAS numbers for chemicals. Where applicable, the
encoding adds the external identifier as a separate triple, so a
later cross-corpus alignment pass can merge candidate splits via
\texttt{owl:sameAs} against the Wikidata IRI even if the local
IRIs were minted independently.

\subsection{Inclusion criteria for the controlled vocabulary}

\begin{itemize}
\item \textbf{Belongs in \texttt{mlips-vocab.ttl}:}
  recurs across multiple MLIP papers, has stable identity outside
  any paper, will recur as the corpus grows, has or could have an
  external identifier (Wikidata, DOI). Examples: XC functionals
  (\texttt{mlips:PBE}, \texttt{mlips:HSE06}); pseudopotential
  types (\texttt{mlips:PAW}); wave-function methods
  (\texttt{mlips:CCSDT}); MLIP method families
  (\texttt{mlips:MTP}, \texttt{mlips:MACEMethod}); software
  (\texttt{mlips:VASP}, \texttt{mlips:MLIP2}, \texttt{mlips:FHIaims}).

\item \textbf{Stays paper-local:}
  specific MLIP runs, trained models, hyperparameter settings,
  training datasets that are themselves paper-specific (the
  dataset \emph{instance}; if it's a published reusable dataset,
  link via $\dlRole{sameAsWikidata}$ or a DOI to the canonical
  one), specific atomic configurations, specific benchmark
  results.

\item \textbf{Borderline cases.}
  A reusable published dataset (e.g., MPtrj) gets a vocabulary
  IRI; each paper's specific subset stays paper-local with
  \texttt{prov:wasDerivedFrom} pointing at the canonical one. A
  software version (\texttt{MLIP-2 v2.0.1}) is the software in
  the vocabulary plus a per-paper \texttt{Implementation} node
  with the version as a literal.
\end{itemize}

\subsection{Workflow when a new canonical entity is encountered}

The agent must not be blocked by the vocabulary, and the vocabulary
must not grow uncontrolled.

\begin{enumerate}
\item \textbf{Encode locally as a candidate.} The agent mints a
  paper-local IRI and tags it with the new
  \texttt{mlips:candidateForVocabulary} object property whose
  range is the appropriate vocabulary class. Example:
\begin{lstlisting}[language=turtle]
ex:xc-rev-pbe-d3-musaelian2023
    a mlips:XCFunctional ;
    rdfs:label "revPBE-D3" ;
    mlips:candidateForVocabulary mlips:XCFunctional .
\end{lstlisting}
  The encoding still passes round-trip; the candidate flag is a
  separate triple.

\item \textbf{Audit pass} runs after every encoding batch and
  collects all $\dlRole{candidateForVocabulary}$ triples across
  the corpus. Output: a review queue.

\item \textbf{Curator decides} per candidate, applying the
  inclusion criteria above:
  \begin{itemize}
  \item \emph{Promote.} Add the IRI to \texttt{mlips-vocab.ttl}
    with full metadata (label, comment, sameAsWikidata where
    applicable).
  \item \emph{Merge.} The candidate is the same as an existing
    canonical entity, modulo the local label; assert
    \texttt{owl:sameAs} or rewrite the IRI in the paper's .ttl.
  \item \emph{Keep paper-local.} The candidate is genuinely
    paper-specific and was wrongly tagged; remove the
    candidate flag.
  \end{itemize}

\item \textbf{Re-encoding pass} (mechanical) rewrites paper-local
  IRIs to canonical ones if (a) or (b) was chosen. Round-trip
  check confirms no semantic drift.
\end{enumerate}

\subsection{Audit queries (sketch)}

Cross-corpus duplicate-detection queries to run as part of the
review:

\begin{itemize}
\item Same \texttt{rdfs:label}, same \texttt{rdf:type}, different
  IRI: likely split.
\item Same \texttt{mlips:chemicalFormula}, same
  \texttt{mlips:materialClass}, different IRI: likely split (a
  material).
\item Same \texttt{rdfs:label} on a $\LibraryC$ instance,
  different IRI: software-package split.
\item Same IRI with conflicting \texttt{rdf:type} or
  contradictory data property values across paper graphs:
  collision.
\end{itemize}

These are 5--10-line SPARQL queries; we add them to
\texttt{artifacts/kg/audit/} when the controlled vocabulary
infrastructure lands.

\subsection{Concon implications}

When the corpus moves to Concon, each paper occupies a named graph
\texttt{<concon://graphs/\textit{paper-id}>}. Cross-graph queries
treat IRIs as graph-local unless explicitly merged. The MCP
\texttt{submit\_encoding} tool runs the audit server-side and
rejects submissions that mint a new IRI for an entity already in
the controlled vocabulary, returning a hint like ``You used
\texttt{ex:pbe-functional-foo2024}; the canonical IRI is
\texttt{mlips:PBE}. Resubmit?''. Vocabulary additions go through a
curator interface (or admin agent privilege).

%% ====================================================================
